% BHU PYQ -- Manually transcribed & error-corrected
\documentclass[11pt,a4paper]{article}

\usepackage[a4paper,top=2.5cm,bottom=2.5cm,left=2.8cm,right=2.8cm,headheight=14pt]{geometry}
\usepackage{amsmath,amssymb}
\usepackage[T1]{fontenc}
\usepackage[utf8]{inputenc}
\usepackage{lmodern}
\usepackage{microtype}
\usepackage{fancyhdr}
\usepackage{enumitem}
\usepackage{hyperref}
\usepackage{lastpage}
\usepackage{parskip}

\hypersetup{colorlinks=true,urlcolor=black,linkcolor=black,
  pdftitle={BSC-07A: Ancillary Physics-II -- BHU PYQ},pdfauthor={Banaras Hindu University}}

\pagestyle{fancy}\fancyhf{}
\renewcommand{\headrulewidth}{0.4pt}\renewcommand{\footrulewidth}{0.4pt}
\fancyhead[L]{\small   \textit{Banaras Hindu University}}
\fancyhead[R]{\small   \textit{BSC-07A: Ancillary Physics-II}}
\fancyfoot[C]{\small Page  \thepage\ of \pageref{LastPage}}


\newlist{parts}{enumerate}{2}
\setlist[parts,1]{label=(\alph*),leftmargin=*,itemsep=5pt,topsep=3pt}
\setlist[parts,2]{label=(\roman*),leftmargin=*,itemsep=3pt,topsep=2pt}

\newcommand{\pts}[1]{\hfill   \textit{[#1]}}


\begin{document}


\begin{center}
  {\large\bfseries BANARAS HINDU UNIVERSITY}\\[2pt]
  {
\normalsize B.Sc. (Hons.) Semester IV Examination, 2023-24}\\[6pt]
  \rule{0.8\linewidth}{0.5pt}\\[6pt]
  {\large\bfseries Paper No. BSC-07A: Ancillary Physics-II}\\[4pt]
  {
\normalsize Time: 3 Hours\hfill Maximum Marks: 70}
\end{center}
\rule{\linewidth}{0.4pt}
\smallskip



\noindent   \textit{Instructions: Answer   \textbf{five} questions including Question~1 which is compulsory. Symbols have their usual meanings.}
\bigskip




\noindent   \textbf{Question 1.} Answer any   \textbf{five} of the following: \pts{5 $ \times$ 4 = 20}

\begin{parts}
  \item Define the interference and diffraction of light with suitable diagrams.
  \item On which law of thermodynamics is a thermometer based? Explain your answer.
  \item Discuss polarization by reflection with a suitable diagram.
  \item How do you measure specific heat capacity? Why does sand heat up faster than water on beaches?
  \item Explain one biological process that uses a law of thermodynamics.
  \item Define the entropy of a system. In which thermodynamic process does the entropy not change and why?
  \item Write the expression for the focal length and power of a thick lens of thickness $t$ and radii of curvature $R_1$ and $R_2$.
  \item Consider a bacterium of mass $10^{-12}\, \text{g}$. Find the mass of one mole of bacteria in units of elephant masses (assume an average elephant mass of $6000\, \text{kg}$).
  \item Derive the ideal gas equation by using the empirical laws for gases. State one application of the ideal gas equation.
\end{parts}

\medskip


\noindent   \textbf{Question 2.} 

\begin{parts}
  \item Explain the diffraction pattern due to a straight edge with a suitable diagram and write the expressions for maximum and minimum intensities. \pts{6}
  \item Define polarization with a suitable diagram and discuss the difference between longitudinal and transverse waves. \pts{8}
\end{parts}

\medskip


\noindent   \textbf{Question 3.}

\begin{parts}
  \item Discuss the differences between a telescope and a microscope and their respective applications. \pts{4}
  \item What is the refractive index of a medium? The absolute refractive indices of glass and water are $3/2$ and $4/3$, respectively. If the speed of light in glass is $2  \times 10^8\, \text{m/s}$, calculate the speed of light in air and water. \pts{5}
  \item Define Fresnel and Fraunhofer diffraction with suitable diagrams. \pts{5}
\end{parts}

\medskip


\noindent   \textbf{Question 4.}

\begin{parts}
  \item Explain Carnot's cycle using a schematic diagram. A Carnot engine operates between reservoir temperatures of $800\, \text{K}$ and $300\, \text{K}$. Calculate the percentage efficiency of the engine. \pts{9}
  \item What is Brownian motion? Prove that heavier molecules exhibit slower Brownian motion. \pts{5}
\end{parts}

\medskip


\noindent   \textbf{Question 5.}

\begin{parts}
  \item Helmholtz free energy is given by $F = U - TS$, where $U$ is internal energy, $T$ is temperature and $S$ is entropy. Deduce the expressions for entropy and pressure as negative rates of change of the Helmholtz function at fixed volume and temperature, respectively. \pts{6}
  \item State the first law of thermodynamics. Apply it to ideal gas expansion and comment on the result. \pts{6}
\end{parts}

\hfill   \textbf{OR}


\begin{parts}
  \item Show that when unpolarized light passes through a polarizer, the intensity of the transmitted light is exactly half of the incident light intensity. \pts{6}
  \item What is a diffraction grating? If the width of each slit is $b$ and the separation between the two slits is $a$, write the expression for the intensity distribution and draw the diffraction pattern. \pts{6}
\end{parts}

\vfill
\rule{\linewidth}{0.4pt}

\begin{center}\small
  \href{https://bhu-exam.vercel.app/}{Click here to visit our website}\\[4pt]\href{https://me-aryan.vercel.app/}{Click here to visit our contributor}\\[4pt]\href{https://quashan-me.vercel.app/}{Click here to visit our contributor}
\end{center}

\end{document}
